(1) 由时变电偶极的迟滞场远区近似，$\boldsymbol E_{\rm rad}=[\ddot{\boldsymbol p}(t-r/c)\times\hat r]\times\hat r/(4\pi\epsilon_0c^2r)$，对轴向偶极其非零横向分量的幅值为$E_{\theta,0}=p_0\omega^2\sin\theta/(4\pi\epsilon_0c^2r)$。辐射磁场为$H_\phi=E_\theta/Z_0$、$Z_0=\sqrt{\mu_0/\epsilon_0}=\mu_0c$，两场同相且正交，近场$1/r^2,1/r^3$项在此尺度被忽略。
(2) 周期平均Poynting径向分量$\langle S_r\rangle=\langle E_\theta H_\phi\rangle=E_{\theta,0}^2/(2Z_0)=p_0^2\omega^4\sin^2\theta/(32\pi^2\epsilon_0c^3r^2)$。故单位立体角的功率$dP/d\Omega=r^2\langle S_r\rangle$与$\sin^2\theta$成正比，轴向$\theta=0,\pi$没有辐射，而赤道方向最强。
(3) 利用$\int_{\mathbb S^2}\sin^2\theta\,d\Omega=8\pi/3$，得到$\overline P=p_0^2\omega^4/(12\pi\epsilon_0c^3)$。若偶极来自单电荷$q$的小幅运动$p_0=qx_0$，瞬时Larmor功率$q^2a^2/(6\pi\epsilon_0c^3)$，其中$a(t)=-\omega^2x_0\cos\omega t$，对$\cos^2$取平均为$1/2$，恰复现上述$\overline P$。
(4) 因$\dot p(t)$为电流，幅值关系$I_0=\omega p_0/\ell$，设等效天线辐射功率$\overline P=I_0^2R_{\rm rad}/2$，解得$R_{\rm rad}=\omega^2\ell^2/(6\pi\epsilon_0c^3)=Z_0(k\ell)^2/(6\pi)=80\pi^2(\ell/\lambda)^2\,\Omega$。只有$k\ell\ll1$、电流分布能用有效长度$\ell$表示时才适用；真实线天线的输入阻抗还含电抗及导体损耗。
